\documentclass[a4paper]{article}
\usepackage{graphicx}
\usepackage{xcolor}
\usepackage[margin=1in]{geometry}
\usepackage[utf8]{inputenc}
\usepackage[absolute]{textpos}
\usepackage{listings}

\title{

\\
{DevOps}
}

\author{Millennium}

\begin{document}

\begin{titlepage}

\begin{center}

{\Huge \bf BeakPeek}\\[10pt]
\rule{\textwidth}{1pt}\\[20pt]
{\bf \huge \sc Millennium} \\[10pt]
{\huge University Of Pretoria} \\[15pt]
\includegraphics[width=7cm]{../../../res/MillenniumLogo.jpeg}\\[15pt]
{\huge Algorithmic Innovation Award}\\[60pt]
\rule{8cm}{1pt}\\[5pt]
{\bf \huge \sc In Partnership With}\\[15pt]
{\huge \bf AgileBridge}\\[5pt]
\rule{8cm}{1pt}\\[50pt]

\begin{minipage}[t]{10cm}
	{\Large{\bf Group:\\ 22}}
\end{minipage}\hfill\begin{minipage}[t]{5cm}\raggedleft
	{\Large{\bf Mentor: \\Ayaaz\\ }}
\end{minipage} \\[90pt]
{\Large 2024/10/18} \\ [5pt]

\rule{\textwidth}{1pt}\\[10pt]

\end{center}

\vfill

\end{titlepage}

\newpage

\section{Introduction}

We as Millennium have from the start of the capstone project placed a high 
level of importance on DevOps which was managed by and maintained by Ivan Horak
who has been our DevOps engineer.\newline

Seeing that the scale of this project is so large we feel it is important that we 
give a summary of what BeakPeek is and how we have structured the repository and 
all of the different parts of how BeakPeek functions.

\subsection{What is BeakPeek}
BeakPeek is a bird watching app that intends to making birding as easy as fly.\newline

Our core functional requirements are:
\begin{itemize}
    \item Data ingress and transformation
    \item Usage of the data from SABAP (South African Bird Atlas Project) for
        \begin{itemize}
            \item Heat map generation by pentad (A 5 minute by 5 minute area)
            \item Bird population estimation
            \item Bird quiz generation
        \end{itemize}
    \item User management
    \item Bird Achievements
    \item Lifelists (list of birds that a user has seen)
\end{itemize}

\subsection{Project/Repository structure}

Our project has been split into a couple common directories. 
Those being:
\begin{itemize}
    \item doc: For documentation
    \item res: For resources
    \item scripts: For bash and SQL scripts
    \item beakpeek: For our flutter frontend
    \item dotnet: For our two .Net backends
        \begin{itemize}
            \item BirdApi: For the bird API 
            \item UserApi: For our user portal and management API 
        \end{itemize}
\end{itemize}

\subsection{What we use?}
The tools that we have used are:
\begin{itemize}
    \item Azure to host the bird API and User API
    \item .Net entity framework core for both API's
    \item .Net identity for the User portal and API
    \item Flutter for mobile app development
    \item MSSQL for both databases also on Azure
\end{itemize}

\subsection{Most Important things to note for DevOps}

\begin{itemize}
    \item We use Azure to host our 2 .Net backends/API's
    \item We use MSSQL for our database
    \item We have an Android release on Google Play
    \item We use GitHub actions for our CI/CD Workflows
    \item We get our data from SABAP
\end{itemize}

\section{Non-Functional Requirements}

We as Millennium set out for ourselve 3 main Non-Functional requirements that we
thought would be the most important and applicable to BeakPeek and would best show 
our wide range of skills and show our abilities as software engineers.

\subsection{Performance}
One of the most important skills as software engineers is the ability to make performant
and scalable solutions. \newline

\subsubsection{Load Testing}
We use Azure Load Testing to simulate 50 users constantly signing up, logging in, 
getting their profile, updating their profile and then finally deleting their profile

To show that we can make highly performant and highly scalable solution we have 
used a load test which is hosted on azure in conjunction with our api's. \newline

To ensure that our users are never unable to login and manage their profiles we 
used a load test to ensure that even if and when we have high volumes of users we 
never drop a request. This is done by simulating 50 users simultaneously and 
constantly signing up, logging in, getting their profile, updating their profile
and finally deleting their account. This is done over 10 minutes which gives us
approximately 2.49 thousand requests with an error percentage of 0.36\%. This 
shows that our User API can handle a high qauntity of users easily.


\subsubsection{Statistics}

\subsection{Security}
To ensure our users security and ensure that their data stays secure and that 
only users who are authenticated are allowed to login and edit their data we 
have used Json Web Tokens which are an easy to transport and easy to use authentication
token as it ensures that only users who have been issued a token from our user Api
as anyone could decrypt a JWT but only JWTs issued by our user api are valid as they 
are validated against a secret key on our user API.

\subsubsection{User Portal/Api}

We used the .Net Identity framework for our user management because it is a widely 
used and trusted tool made by Microsoft which meant that it worked smoothly with 
the tools built into Azure and enabled a much more comprehensive DevOps experience.
\newline
Using Identity framework also allowed for easy integration of other login and web auth 
providers such as Google and also allowed easy role based authorization which was important 
as the user portal/api also serves as an admin portal allowing us to easily give manage 
our user, roles and Achievements which were also integrated into the user api along with 
the users life list and other user claims.

\subsubsection{Usability}


\newpage

\section{CI/CD tools and pipelines}

\subsection{Tools}

To ensure the best possible developer experience we used a multitude of tools and integrations
in the CI/CD pipelines. \newline

Some of the most important tools were:
\begin{itemize}
    \item Fastlane for easy builds and releases to Google play for our mobile app
    \item JMeter was used for load testing our Api
    \item Codecov for unit testing tracking
    \item Github actions for automated workflows and tasks
\end{itemize}

\subsubsection{Fastlane}

Fastlane is a tool made especially for the building and releasing of mobile apps to the Google 
Play store and to Apples App store. It is configured in ruby and has plenty of integrations,
such as for deploying using firebase and of course building different flutter environments.


\subsubsection{JMeter}

JMeter was used primarily for our perfomance requirement as JMeter allows for easy configuration
using different inputs such as CSV files which as useful for ensuring that every user that 
was created in the load test was a unique user and was consistent across different runs.

\subsubsection{Codecov}
To ensure that our the anything we make is up to the proper standards and that 
there are no unintended consequences of any new features we add we have employed
the use of Codecov to be able to track our code coverage across our project and 
across the different sections of the project.
\newline

We have split up our code coverage tracking into two primary areas, the flutter frontend and the 
.Net backend. This allowed us to easily track and monitor changes in the environment and ensure
that any time we made a release or deployment we could be certain that there would be minimal 
issues in production.

To ensure that our code is properly tested and up to industry standards we have
aimed for atleast a 75\%  line coverage for each different scope of our project.
\newline
In this regard we have managed to maintain a 91\% code coverage for the BirdApi.

\subsection{Workflows}

Probably the most important and useful tool at our disposal was GitHub actions as it allowed for 
easily configurable and integrated automations in the form of workflows for automating tedious or 
otherwise impossible tasks directly from GitHub.\newline

All of the workflows were written by our DevOps engineer (Ivan Horak), who ensured that our developer
experience was always as easy as it possibly could be. 
\newline 

Ivan wrote a number of workflows and full pipelines to automate most of the tasks that are often 
done manually. Bellow is a list of some of the workflows and pipelines. 

\subsubsection{Discord Notify}

The discord notification workflow is a reusable workflow used to send embed messages 
to the Millennium discord server which allows for easy observability of what is happening
in the project and if there are any failing workflows, releases or deployments.

\subsubsection{Deploy BirdApi, UserApi}

There are two main workflows for deploying the BirdApi and UserApi to Azure, these are the 
essential for the easy management of the two backends as it gave us a lot of control over 
how and when the APIs would be deployed.
\newline \newline

These workflows also use best practices in that they cache any dependencies for future workflows 
which significantly reduces deployment times.

\subsubsection{Update Bird data}

The update bird data workflow is a custom workflow that uses in house scripts to download 
all 9 provinces data as CSV's along with a CSV for all the birds in South Africa which it 
then uses to import all the data into the Bird database using the Bird API which 
has a couple endpoints that accept CSV files as input and transform the data in the files to match
our Bird and Province models. These CSV files have approximately 1.5 million rows of data 
that would be difficult for anyone to manage efficiently. For our use case some of the 
birds that would get fetched from SABAP would have invalid data such as a low report rate 
which don't work with the population estimates that we make using the data that is imported 
through this workflow. 

\subsubsection{Label Pull Requests}

The label pull requests workflow is a simple yet useful workflow that is used to add the appropriate
labels to any pull requests based on the branch name and contents which are configured 
in the labeler.yml file making for easy extensibility and ensuring that every pull
request is always appropriately labeled. 

\subsection{The coverage workflows}

The coverage workflows are run for the flutter frontend and the .Net backend. These are usually run 
together as Codecov does not track any directories that are not included in the commit. This proves 
problematic as more often than not the frontend and backend workflows are not run together, as 
changes should very rarely effect both sources in the same commit. To get around this problem 
we have made use of artifacts that are used to store coverage reports in between workflows for later
useage, Ivan has leveraged artifacts by first checking if there is an existing artifact with the 
same name as the hash of the different sources package lock files, this allows us to only run 
the tests when a change that would effect the test runs output. We have used this to in a sense 
short circuit the coverage uploads as if there is an artifact for either of the reports that will be 
used alongside any new reports that are generated by these workflows. 
\newline \newline

These workflows can still run independently if needed allowing for rapid iteration. We have also 
used a seperate workflow as a sort of manager for both of these workflows that if they need to be 
run without changing the actual source code, the workflows can still be run either manually or 
automatically.

\subsubsection{Publish Android release}

The publish Android release is a workflow that uses the Fastlane tool set to build, test and release 
the flutter frontend as an Android aab file which is a requirement of Google for more secure 
applications. This workflow also uses dependency caching to allow for faster runs of the workflows.

\subsection{Pipelines}

To connect all of the seperate workflows that have been mentioned above there 
is a full CICD pipeline that is used to connect all of the other workflows 
in a way that ensures that all of the industry best practices are used. This
is done by chaining all of the workflows together to ensure that if any step 
in the build, test or deployment stages fails nothing permanent is 
done to either of the production databases and the Google play release along with the 
deployed APIs. 
\newline \newline 

The full pipeline is made up of all of the other workflows in such a way that discord notifications 
are sent at each stage of the pipeline with links to any build artifacts or errors from any of 
the steps and or jobs. The final steps in the pipeline is releasing the app and also updating the 
bird data.

\newpage
\section{Observability Principles}

To ensure that we follow all of the best observability principles Ivan has put together a 
couple different tools for monitoring the APIs and database as well as the traffic 
to the APIs. He has also added in other methods which are used to monitor workflows and 
builds as well as the projects repository itself.

\subsection{Dashboards}

One of the primary tools that we have used to monitor our APIs is through the use of 
different Dashboards for all of the different tools that we use. 

\subsubsection{Azure Metrics}

Azure offers an extensive list of tools that can be used and that we use to make dashboards for 
monitoring metrics like the response time of the APIs or the usage of the CPU for the 
different virtual machines that are hosted on Azure. 

\subsection{Azure Monitoring}

Azure also offers a service called application insights which we have used for tracking the hosted APIs on Azure and sending alerts if an error or some exception is raised. 
\newline \newline

Azure also offers tools for automatic up scaling and downscaling of the resources that we manage 
and host on Azure. We have in particular used the health check service to ping the server and 
alert us if any of the containers become unhealthy.

\subsection{Azure Logs}

We have also used the azure storage account along with the provided logging tools to ensure that 
if any of our deployments go down or stop working for any reason we can always go back and 
look at what happened at the time of the incident.

\subsection{GitHub Actions}
We also use GitHub actions to ensure that we have maximum observability at any given time. As 
we can't always be at our computers however all of us do have discord on our phones which 
allows us to get notified easily and quickly any time an incident occurs.

\newpage
\section{Operational Practices}

In terms of the operational practices we have used to ensure that we can always recover from 
a fatal mistake, we have made backups of the data stored on SABAP to ensure that we its 
unlikely that we could find ourselves in a situation where we cannot recover any lost bird data. 
\newline \newline

The bird data is also automatically updated once a day through the usage of the update bird 
date workflow which has a schedule that it runs on as well as any other actions that can 
cause the workflow to be run, such as the BirdApi being deployed.

\subsection{Data Ingress}

As it was mentioned above we ensure that any data we use is easily managed and should never require 
manual transformations or editing. This extends to the bird images which are stored in an Azure 
storage account and are automatically fetched and updated as bird info is requested.

\subsection{Users data is stored locally on device}

One of the features that we implemented for BeakPeek was the ability for the mobile app to run 
without an internet connection. This means that even if a problem were to arise on our side our 
users should be minimally effected.

\subsection{Repository Management}

Finally the most important resource we have is the repository and to make sure that nothing 
effects the operation and usage of the repository we have put Ivan in charge of anything related 
to it. In this role Ivan reviews all pull requests ensuring that they use proper form and 
proper descriptions and that there are no other problems.

\end{document}
